% ID: SEL-D0072-OB-E03
% ESTADO: validada. S072.
\textbf{Parte I.} Sea $E$ «es español» y $M$ «es menor de edad».

\textbf{a)} Son mayores de edad el $49\%$ del total que además es español y el $30\%\cdot60\%=18\%$ que no es español. Sumando ambos grupos se obtiene $\boxed{0{,}67}$.

\textbf{b)} Entre los menores, cuya probabilidad total es $0{,}33$, los que no son españoles tienen probabilidad conjunta $0{,}30\cdot0{,}40=0{,}12$. Entonces $P(E^c\mid M)=0{,}12/0{,}33=\boxed{4/11}$.

\textbf{Parte II. a)} El error estándar es $2/\sqrt{64}=0{,}25$ horas. Al $97\%$ usamos $z_{0{,}985}\approx2{,}17009$. El intervalo normal es
\[
7\pm2{,}17009\cdot0{,}25
\approx\boxed{[6{,}4575,\ 7{,}5425]}\ \text{horas}.
\]
\textbf{b)} Imponemos $z_{0{,}985}\cdot2/\sqrt n\le0{,}25$. Resulta $n\ge(8z_{0{,}985})^2\approx301{,}39$, de modo que el mínimo es $\boxed{302\text{ jóvenes}}$.
